%% 5-12 物理での積分：仕事と流量（第11回：ページ制限なし。積分の3段「分ける・かける・たす」を、量ごとに図と数で見せる）
\bkchapter{5}{12}{物理での積分：仕事と流量}{\bklevel{3}}{物理での積分：仕事と流量}
\bkmichishirube{2-15,4-12,5-11}{5-13,6-15}

\begin{goal}
\begin{itemize}[label={\color{bkTheme}●}, leftmargin=1.6\zw, itemsep=0pt, topsep=1pt]
\item \bklevel{3}\enspace 力が変わるときの仕事を、$F$-$x$ グラフの面積（積分）で求められる。
\item \bklevel{3}\enspace 流量や電力を時間で積分して、たまった量を求められる。
\item \bklevel{3}\enspace 積分を使って、平均の値を求められる。
\end{itemize}
\end{goal}

\begin{hatena}
\noindent\begin{minipage}[c]{0.58\linewidth}\raggedright
ばねをのばすとき、はじめは軽い力でのびますが、のばすほど強い力が必要になります。ばねの力は、のびに比例するからです（2-15）。

ばねを\qty{10}{cm}のばすのに必要な「仕事」は、いくらでしょう。仕事は「力×距離」で求めます。でも、力が変わり続けるので、どの力を使えばよいかわかりません。
\end{minipage}\hfill
\begin{minipage}[c]{0.38\linewidth}
\begin{tikzpicture}[x=1mm, y=1mm]
  \fill[bkGray!40] (-2,-6) rectangle (0,6);
  \draw[bkInk, line width=0.8pt, decorate, decoration={coil, aspect=0.5, segment length=2.4mm, amplitude=2.2mm}] (0,0) -- (26,0);
  \fill[bkTheme!40] (26,-4) rectangle (32,4);
  \draw[bkInk, line width=0.6pt] (26,-4) rectangle (32,4);
  \draw[bk arrow, bkOrange] (33,0) -- (44,0);
  \node[font=\sffamily\scriptsize, anchor=south] at (39,1) {$F=kx$};
  \draw[bkGray, <->, line width=0.5pt] (18,-8) -- (26,-8);
  \node[font=\sffamily\scriptsize, anchor=north] at (22,-9) {のび $x$};
\end{tikzpicture}
\end{minipage}
\end{hatena}

5-10と5-11で学んだ積分は、「細かく分けて、かけて、たす」計算でした。速度から距離を求めるだけでなく、物理のいろいろな「たまる量」を計算する道具です。この章では、仕事、水の量、電気のエネルギー、電気の量を、この3つの段で求めます。どの量でも、同じ3つの段をくり返すだけです。

\clearpage
%% ---------------- やってみよう・思い出そう ----------------
\begin{yattemiyou}[ばねばかりで、のばす力をはかろう]
\textbf{用意するもの}\par
ばねばかり（\qty{20}{N}まで）、ばね、ものさし
\begin{enumerate}[label={\bkcircled[bkGreen]{\arabic*}}, leftmargin=2.4\zw, itemsep=1pt, topsep=3pt]
\item ばねの一方を固定し、もう一方をばねばかりで引く。のびを\qty{2}{cm}ずつふやしながら、力を読もう。
\item 横軸をのび、たて軸を力にしてグラフをかこう（$F$-$x$ グラフ）。
\end{enumerate}
\begin{bkcenter}
\begin{bktable}{colspec={Q[l,m,wd=30mm] *{6}{Q[c,m,wd=10mm]}}, rowsep=1.5pt}
のび $x$（cm） & 0 & 2 & 4 & 6 & 8 & 10 \\
力 $F$（N） & 0 & 4 & 8 & 12 & 16 & 20 \\
\end{bktable}
\bktabcaption{ある班の結果。力はのびに比例している。}
\end{bkcenter}
\end{yattemiyou}

\noindent \qty{2}{cm}のびるごとに、力は\qty{4}{N}ずつふえています。のびを m（メートル）で表すと、\qty{2}{cm}は
\[ \qty{2}{cm}\div\qty{100}{cm/m}=\qty{0.02}{m} \]
です（\qty{100}{cm/m} は「\qty{1}{m}あたり\qty{100}{cm}」。2-10）。だから、のび\qty{1}{m}あたりの力（ばね定数 $k$）は
\[ k=\qty{4}{N}\div\qty{0.02}{m}=\qty{200}{N/m} \]
です。このばねの力は $F=200x$（$x$ は m、$F$ は N）と書けます。たとえば $x=0.1$ なら $F=200\times0.1=20$ で、表の\qty{20}{N}と合います。\bkidx{ばねていすう}{ばね定数}

\begin{omoidasou}
\bkrecall{4-12}{物を力の向きに動かしたとき、仕事 $W$＝力 $F$×動いた距離 $s$。単位は J（ジュール）。\qty{1}{N}の力で\qty{1}{m}動かしたときの仕事が\qty{1}{J}（1 J＝1 N·m）。}
\bkrecall{2-15}{ばねの力は、のびに比例する：$F=kx$。$k$ はばね定数（N/m）。}
\bkrecall{5-10}{定積分 $\displaystyle\int_a^b f(x)\,\dd x$ は、高さ $f(x)$、はば $\Delta x$ の細い長方形を $a$ から $b$ までたし、はばを0に近づけた先。グラフの下の面積。}
\bkrecall{5-11}{$\displaystyle\int_a^b f(x)\,\dd x=\Bigl[\,F(x)\,\Bigr]_a^b=F(b)-F(a)$（$F$ は微分すると $f$ になる関数）。$\displaystyle\int x\,\dd x=\frac{x^2}{2}+C$}
\bkrecall{2-10}{換算は「1をかける」：\qty{60}{s/min}、\qty{3600}{s/h}、\qty{1000}{W/kW}。単位をそろえてから計算する。}
\end{omoidasou}

%% ---------------- 数学の道具 ----------------
\bksection{doguu}{数学の道具}

\bksubsection{1}{力が一定なら、仕事は長方形の面積}

\label{bk:5-12:shigoto}物を動かすたいへんさは、力が大きいほど、動かす距離が長いほど大きくなります。たとえば、重い荷物を2つ運ぶのは1つのときの2倍たいへんで、2倍の距離を運ぶのも2倍たいへんです。そこで物理では、
\[ \text{仕事}=\text{力}\times\text{力の向きに動いた距離},\qquad W=Fx \]
と決めています（4-12）。力が\qty{1}{N}、距離が\qty{1}{m}なら\qty{1}{J}です。これを\textbf{仕事}といいます。\bkidx{しごと}{仕事}

\noindent\begin{minipage}[c]{0.46\linewidth}
\begin{tikzpicture}
\begin{axis}[bkplot, width=56mm, height=40mm, xmin=0, xmax=2, ymin=0, ymax=26, xtick={1.5}, ytick={19.6}, xlabel={$x$（m）}, ylabel={$F$（N）}]
  \fill[bkOrange!25] (axis cs:0,0) rectangle (axis cs:1.5,19.6);
  \addplot[bkcurve, line width=1.4pt] coordinates {(0,19.6) (1.5,19.6)};
  \node[font=\sffamily\small, fill=bkOrange!25, inner sep=1pt] at (axis cs:0.75,9) {\qty{29.4}{J}};
\end{axis}
\end{tikzpicture}
\end{minipage}\hfill
\begin{minipage}[c]{0.52\linewidth}\raggedright
横軸を位置 $x$、たて軸を力 $F$ にした $F$-$x$ グラフでは、力が一定なら、仕事はたて $F$、横 $x$ の\textbf{長方形の面積}です。$v$-$t$ グラフの長方形の面積が距離になったのと同じ形です（5-9）。単位も N×m＝J になります。
\end{minipage}

\noindent\textbf{例}\quad 質量\qty{2}{kg}の荷物を、ゆっくり一定の速さで\qty{1.5}{m}持ち上げます。地球は、質量\qty{1}{kg}あたり約\qty{9.8}{N}の力（重力）で、物を下に引いています。この9.8は、5-5で調べた落ちる物の加速度 $g=\qty{9.8}{m/s^2}$ と同じ数です。だから荷物にはたらく重力は
\[ F=mg=\qty{2}{kg}\times\qty{9.8}{N/kg}=\qty{19.6}{N} \]
です。ゆっくり一定の速さで持ち上げるときは、この重力と同じ大きさの力で上に引けばよいので、力は\qty{19.6}{N}で一定です。
\[ W=Fx=19.6\times1.5=\qty{29.4}{J} \]
上の図の長方形の面積です。

\bksubsection{2}{力が変わるときは、細かく分けて、かけて、たす}

\label{bk:5-12:sekibun}\bkyattemiyouref のばね（$F=200x$）を、0から\qty{0.1}{m}までのばす仕事を考えます。力が変わるので、「力×距離」をそのまま使えません。そこで、5-10と同じ3つの段で考えます。

\noindent\textbf{段1　細かく分ける}\quad 0から\qty{0.1}{m}までを、5つに等しく分けます。1つのはばは $0.1\div5=\qty{0.02}{m}$ です（長さ0.1を5こに等しく分けるから）。

\noindent\textbf{段2　かける}\quad 細い区間の中では、力はほとんど変わりません。そこで、区間の右はしの力を、その区間の力とみなします。たとえば2番めの区間（0.02から\qty{0.04}{m}）の右はしは $x=0.04$ なので、力は $200\times0.04=\qty{8}{N}$。この区間の小さな仕事は、力×はばで $8\times0.02=\qty{0.16}{J}$ です。

\noindent\textbf{段3　たす}\quad 5つの小さな仕事をたします。

\begin{bkcenter}
\begin{bktable}{colspec={Q[l,m,wd=32mm] *{5}{Q[c,m,wd=12mm]} Q[c,m,wd=12mm]}, rowsep=1.5pt, colsep=4pt, row{1}={bg=bkTheme!10}}
区間の番号 & 1 & 2 & 3 & 4 & 5 & 合計 \\
右はしの位置（m） & 0.02 & 0.04 & 0.06 & 0.08 & 0.1 & \\
右はしの力（N） & 4 & 8 & 12 & 16 & 20 & \\
小さな仕事（J） & 0.08 & 0.16 & 0.24 & 0.32 & 0.40 & 1.2 \\
\end{bktable}
\bktabcaption{右はしの力を使った和。小さな仕事＝右はしの力×0.02。}
\end{bkcenter}

\begin{bkfigure}
\begin{tikzpicture}
\begin{axis}[bkplot, name=s1, width=60mm, height=44mm, xmin=0, xmax=0.115, ymin=0, ymax=24, xtick={0.02,0.04,0.06,0.08,0.1}, xticklabels={,0.04,,0.08,}, ytick={10,20}, xlabel={$x$（m）}, ylabel={$F$（N）}, title={\sffamily\small 右はしの力（大きめ）}, scaled x ticks=false]
  \foreach \k in {1,...,5} {\edef\tmp{\noexpand\fill[bkOrange!30, draw=bkOrange!70, line width=0.5pt] (axis cs:{(\k-1)*0.02},0) rectangle (axis cs:{\k*0.02},{4*\k});}\tmp}
  \addplot[bkcurve, domain=0:0.11] {200*x};
\end{axis}
\begin{axis}[bkplot, at={(s1.east)}, anchor=west, xshift=14mm, width=60mm, height=44mm, xmin=0, xmax=0.115, ymin=0, ymax=24, xtick={0.02,0.04,0.06,0.08,0.1}, xticklabels={,0.04,,0.08,}, ytick={10,20}, xlabel={$x$（m）}, ylabel={$F$（N）}, title={\sffamily\small 左はしの力（小さめ）}, scaled x ticks=false]
  \foreach \k in {1,...,5} {\edef\tmp{\noexpand\fill[bkBlue!25, draw=bkBlue!70, line width=0.5pt] (axis cs:{(\k-1)*0.02},0) rectangle (axis cs:{\k*0.02},{4*(\k-1)});}\tmp}
  \addplot[bkcurve, domain=0:0.11] {200*x};
\end{axis}
\end{tikzpicture}
\bkfigcaption{長方形の面積が小さな仕事。和は、右はしの力で1.2 J、左はしで0.8 J。}
\end{bkfigure}

右はしの力は、その区間でいちばん大きい力なので、1.2 J は少し大きすぎます。左はしの力（いちばん小さい力：0、4、8、12、16 N）を使うと、$0+0.08+0.16+0.24+0.32=\qty{0.8}{J}$ で、少し小さすぎます。ほんとうの仕事は、そのあいだです。分け方を細かくしていくと、次のようになります。

\begin{bkcenter}
\begin{bktable}{colspec={Q[l,m,wd=34mm] *{4}{Q[c,m,wd=17mm]}}, rowsep=1.5pt}
分ける数 & 5 & 10 & 100 & 1000 \\
はば（m） & 0.02 & 0.01 & 0.001 & 0.0001 \\
右はしの和（J） & 1.2 & 1.1 & 1.01 & 1.001 \\
左はしの和（J） & 0.8 & 0.9 & 0.99 & 0.999 \\
\end{bktable}
\bktabcaption{細かく分けるほど、どちらの和も1に近づく。}
\end{bkcenter}

2つの和は、上と下から1にはさまれて近づいていきます。だから、このばねを\qty{0.1}{m}のばす仕事は、ちょうど\qty{1}{J}です。速度から距離を求めたとき（5-10）と、まったく同じ考え方です。

\bksubsection{3}{記号で書くと $\displaystyle\int F\,\dd x$}

\label{bk:5-12:kigou}②の「細かく分けて、かけて、たして、はばを0に近づけた先」を、5-10の定積分の記号で書きます。

\[ W=\int_a^b F(x)\,\dd x \]
$F$-$x$ グラフの、$a$ から $b$ までの面積が仕事です。記号の1つ1つの意味は、下の表のとおりです。

\noindent\textbf{なぜ $\dd x$ を書くの？}\quad 仕事は「力×距離」でした。$F(x)$ が力で、$\dd x$ が（とても小さな）距離です。$\dd x$ を書かないと、かけ算の相手の「距離」がなくなってしまいます。単位も、$F(x)$ の N に $\dd x$ の m がかかって、N×m＝J になります。$\dd x$ は「何をはばにして細かく分けたか」も表しています。ここでは位置 $x$ で分けたので $\dd x$ です。

\begin{bkcenter}
\begin{bktable}{colspec={Q[c,m,wd=26mm] Q[l,m,wd=84mm]}, rowsep=2pt, row{1}={bg=bkTheme!10}}
記号 & 意味 \\
$\displaystyle\int$ & たし合わせる（「和」の sum の S を、たてにのばした形） \\
下の $a$、上の $b$ & はじめの位置（下端）と、おわりの位置（上端） \\
$F(x)$ & 位置 $x$ での力（細い長方形の高さ） \\
$\dd x$ & 小さな距離（細い長方形のはば） \\
$F(x)\,\dd x$ & 小さな仕事（力×小さな距離＝細い長方形の面積） \\
\end{bktable}
\bktabcaption{$\int$ は「インテグラル」と読む（5-10）。}
\end{bkcenter}

\bksubsection{4}{ばねの仕事は $\dfrac{1}{2}kx^2$}

\label{bk:5-12:bane}ばね定数 $k$ のばねを、0から $x$ までのばす仕事を、3つの方法で求めます。どれも同じ答えになります。

\noindent\begin{minipage}[c]{0.44\linewidth}
\begin{tikzpicture}
\begin{axis}[bkplot, width=54mm, height=42mm, xmin=0, xmax=0.125, ymin=0, ymax=25, xtick={0.1}, xticklabels={$x$}, ytick={10,20}, yticklabels={$\frac{kx}{2}$,$kx$}, xlabel={のび}, ylabel={$F$}, scaled x ticks=false]
  \addplot[fill=bkOrange!25, draw=none, domain=0:0.1] {200*x} \closedcycle;
  \addplot[bkcurve, domain=0:0.115] {200*x};
  \draw[bkGreen!70!black, line width=1pt, dashed] (axis cs:0,10) -- (axis cs:0.1,10);
  \node[font=\sffamily\small] at (axis cs:0.072,4.5) {$W$};
  \draw[bkGray, dashed, line width=0.5pt] (axis cs:0.1,0) -- (axis cs:0.1,20);
\end{axis}
\end{tikzpicture}
\end{minipage}\hfill
\begin{minipage}[c]{0.54\linewidth}\raggedright
\textbf{方法1　三角形の面積}\par
$F=kx$ のグラフはまっすぐなので、面積は底辺 $x$、高さ $kx$ の三角形です。
\[ W=\frac{1}{2}\times x\times kx=\frac{1}{2}kx^2 \]
$\dfrac{1}{2}$ は、三角形の面積の「÷2」から来ています。
\end{minipage}

\noindent\textbf{方法2　積分}\quad 力 $ks$（$s$ はとちゅうののび）を、0から $x$ まで積分します。$s$ は、0から $x$ まで動くとちゅうののびです。上端の $x$ と区別するため、別の文字にしています。
\[ W=\int_0^x ks\,\dd s=\Bigl[\,k\times\frac{s^2}{2}\,\Bigr]_0^x=k\times\frac{x^2}{2}-0=\frac{1}{2}kx^2 \]
$\dfrac{1}{2}$ は、こんどは「$s$ を積分すると $\dfrac{s^2}{2}$」（5-11：指数を1ふやして、その数でわる）から出てきました。

\noindent\textbf{方法3　平均の力}\quad 力は0から $kx$ まで、のびに比例して一様にふえます。だから、平均の力は $(0+kx)\div2=\dfrac{kx}{2}$（グラフの緑の点線）です。平均の力で $x$ だけ動かしたと考えると、$\dfrac{kx}{2}\times x=\dfrac{1}{2}kx^2$。三角形を、同じ面積の長方形にならしたことになります。

\noindent\textbf{数で確かめる}\quad \bkyattemiyouref のばね（$k=200$）を $x=0.1$ だけのばすと、
\[ W=\frac{1}{2}\times200\times0.1^2=\frac{1}{2}\times200\times0.01=\frac{1}{2}\times2=\qty{1}{J} \]
②の表で近づいていった1と、ぴったり合います。これが\bkhatenaref の答えです。

\noindent\textbf{まちがえた考え方}\quad 最後の力\qty{20}{N}を使って $20\times0.1=\qty{2}{J}$ とすると、ほんとうの2倍になります。はじめの力0を使えば0 J です。力が変わるときに「どこか1つの力×距離」とすると、まちがえるのです。

のばしたばねは、この仕事の分のエネルギーをたくわえます。手をはなすと、このエネルギーで物をおし出します。これを\textbf{弾性エネルギー}といいます。\bkidx{だんせいえねるぎー}{弾性エネルギー $\frac{1}{2}kx^2$}\bkidx{えねるぎー}{エネルギー}

\bksubsection{5}{流量を積分すると、たまった水の量}

\label{bk:5-12:ryuuryou}じゃ口から1分あたりに出る水の量を\textbf{流量}といい、$Q$（L/min）で表します。L/min は「1分あたり何 L」です。流量が一定なら、たまる水の量は「流量×時間」です。たとえば\qty{10}{L/min}で3分なら、$\qty{10}{L/min}\times\qty{3}{min}=\qty{30}{L}$（min どうしが約分されて L が残る）。速さ×時間＝距離と同じ形です。\bkidx{りゅうりょう}{流量}

\noindent\textbf{流量が変わるとき}\quad ふろのじゃ口を少しずつ開けていくと、流量は $Q=2t$ でふえていきました。$Q$ の単位は L/min、$t$ は開け始めてからの時間（分）です。6分後の流量は、$2\times6=12$ より\qty{12}{L/min}です。6分間で、何 L たまったでしょう。ここでも3つの段を使います。

\begin{enumerate}[label={\bkcircled[bkTheme]{\arabic*}}, leftmargin=2.4\zw, itemsep=1pt, topsep=2pt]
\item \textbf{分ける}：0分から6分までを、1分ずつ6つに分ける。
\item \textbf{かける}：1分のあいだは流量がほぼ一定とみて、右はしの流量×1分。1番めは $2\times1\times1=\qty{2}{L}$、2番めは $2\times2\times1=\qty{4}{L}$、…、6番めは\qty{12}{L}。
\item \textbf{たす}：$2+4+6+8+10+12=\qty{42}{L}$（右はし。大きめ）。左はしの流量では $0+2+4+6+8+10=\qty{30}{L}$（小さめ）。
\end{enumerate}

\noindent\begin{minipage}[c]{0.46\linewidth}
\begin{tikzpicture}
\begin{axis}[bkplot, width=50mm, height=42mm, xmin=0, xmax=7, ymin=0, ymax=14.5, xtick={1,2,...,6}, ytick={4,8,12}, xlabel={$t$（min）}, ylabel={$Q$（L/min）}]
  \foreach \k in {1,...,6} {\edef\tmp{\noexpand\fill[bkBlue!18, draw=bkBlue!60, line width=0.4pt] (axis cs:\k-1,0) rectangle (axis cs:\k,{2*\k});}\tmp}
  \addplot[fill=bkBlue!35, draw=none, domain=0:6] {2*x} \closedcycle;
  \addplot[bkcurve, domain=0:6.6] {2*x};
  \node[font=\sffamily\small, fill=bkBlue!35, inner sep=1pt] at (axis cs:4.3,3) {\qty{36}{L}};
\end{axis}
\end{tikzpicture}
\end{minipage}\hfill
\begin{minipage}[c]{0.52\linewidth}\raggedright
ほんとうの水の量は、42 L と 30 L のあいだです。はばを0に近づけた先は、$Q$-$t$ グラフの下の面積（こい青の三角形）で、
\[ V=\int_0^6 2t\,\dd t=\Bigl[\,t^2\,\Bigr]_0^6=\qty{36}{L} \]
（$t^2$ を微分すると $2t$ なので、$2t$ の原始関数は $t^2$。5-11）三角形の面積でも $\dfrac{1}{2}\times6\times12=36$ で、同じです。
\end{minipage}

\noindent 一般に、時刻 $t_1$ から $t_2$ までにたまる水の量 $V$ は
\[ V=\int_{t_1}^{t_2} Q(t)\,\dd t \]
です。こんどは時間で細かく分けたので、$\dd x$ ではなく $\dd t$（小さな時間）を書きます。$Q(t)\,\dd t$ は「流量×小さな時間＝小さな水の量」です。

\bksubsection{6}{電力を積分すると、電力量}

\label{bk:5-12:denryoku}電気器具が1秒あたりに使うエネルギーを\textbf{電力}といい、単位は W（ワット）です。1秒に\qty{1}{J}使うのが\qty{1}{W}なので、W＝J/s です。たとえば\qty{10}{W}のLED電球は、1秒に\qty{10}{J}のエネルギーを使います。1時間（\qty{3600}{s}）つけておくと、
\[ \qty{10}{J/s}\times\qty{3600}{s}=\qty{36000}{J} \]
です。電力×時間で求めた、使ったエネルギーを\textbf{電力量}といいます。電力が変わるときは、流量のときと同じく、時間で細かく分けてたします。\bkidx{でんりょくりょう}{電力量}
\[ E=\int_{t_1}^{t_2}P(t)\,\dd t \]

\noindent\textbf{kWh（キロワット時）}\quad J では数が大きくなりすぎるので、家庭では電力量の単位に kWh を使います。\qty{1}{kW}で1時間使ったときの電力量が\qty{1}{kWh}です。J に直すには、kW を W に、h を s に換算します。
\[ \qty{1}{kWh}=\qty{1}{kW}\times\qty{1000}{W/kW}\times\qty{1}{h}\times\qty{3600}{s/h}=\qty{3600000}{J}=\qty{3.6e6}{J} \]
（kW と h がそれぞれ約分されて、W×s＝J が残ります。）電力に kW、時間に h（時間）を使って積分すれば、答えは kWh で出てきます。電力に W を使えば、答えは Wh（ワット時。\qty{1}{Wh}は\qty{1}{W}で1時間使った電力量で、\qty{1}{kWh}の $\dfrac{1}{1000}$）です。

\noindent\textbf{例}\quad 電気自動車の充電器の電力は、充電が進むとへっていき、$P=6-t$（kW、$t$ は時間）でした。はじめの2時間の電力量を求めます。$6$ を積分すると $6t$、$t$ を積分すると $\dfrac{t^2}{2}$ なので（5-11）、

\noindent\begin{minipage}[c]{0.44\linewidth}
\begin{tikzpicture}
\begin{axis}[bkplot, width=52mm, height=40mm, xmin=0, xmax=3, ymin=0, ymax=7.5, xtick={1,2}, ytick={4,6}, xlabel={$t$（h）}, ylabel={$P$（kW）}]
  \addplot[fill=bkYellow!45, draw=none, domain=0:2] {6-x} \closedcycle;
  \addplot[bkcurve, domain=0:2.8] {6-x};
  \draw[bkGray, dashed, line width=0.5pt] (axis cs:2,0) -- (axis cs:2,4);
  \node[font=\sffamily\small, fill=bkYellow!45, inner sep=1pt] at (axis cs:1,2) {\qty{10}{kWh}};
\end{axis}
\end{tikzpicture}
\end{minipage}\hfill
\begin{minipage}[c]{0.54\linewidth}\raggedright
\[ \int_0^2(6-t)\,\dd t=\Bigl[\,6t-\frac{t^2}{2}\,\Bigr]_0^2 \]
\[ =\Bigl(12-\frac{4}{2}\Bigr)-0=12-2=\qty{10}{kWh} \]
グラフは $t=0$ で6、$t=2$ で4の直線なので、黄色の台形の面積 $(6+4)\times2\div2=10$ でも確かめられます。
\end{minipage}

\bksubsection{7}{電流を積分すると、流れた電気の量}

\label{bk:5-12:denryuu}電流は、導線のある所を1秒あたりに通りぬける電気の量です。電気の量の単位を C（クーロン）といい、\qty{1}{A}（アンペア）の電流が1秒流れると、\qty{1}{C}の電気が通りぬけます。つまり A＝C/s で、電流は「1秒あたりの電気の量」です。だから、電流 $I$ を時間で積分すると、流れた電気の量 $q$ になります。
\[ q=\int_{t_1}^{t_2}I(t)\,\dd t \]
\noindent\begin{minipage}[c]{0.44\linewidth}
\begin{tikzpicture}
\begin{axis}[bkplot, width=52mm, height=38mm, xmin=0, xmax=13, ymin=0, ymax=0.8, xtick={5,10}, ytick={0.5}, yticklabels={0.5}, xlabel={$t$（s）}, ylabel={$I$（A）}]
  \fill[bkGreen!25] (axis cs:0,0) rectangle (axis cs:10,0.5);
  \addplot[bkcurve, line width=1.4pt] coordinates {(0,0.5) (12,0.5)};
  \draw[bkGray, dashed, line width=0.5pt] (axis cs:10,0) -- (axis cs:10,0.5);
  \node[font=\sffamily\small, fill=bkGreen!25, inner sep=1pt] at (axis cs:5,0.25) {\qty{5}{C}};
\end{axis}
\end{tikzpicture}
\end{minipage}\hfill
\begin{minipage}[c]{0.54\linewidth}\raggedright
\noindent\textbf{例}\quad 電流が\qty{0.5}{A}で一定なら、10秒で
\[ \qty{0.5}{C/s}\times\qty{10}{s}=\qty{5}{C} \]
です。$I$-$t$ グラフでは、たて0.5、横10の長方形（緑）の面積です。電流が変わるときも、$I$-$t$ グラフの下の面積が、流れた電気の量です（\bkreidairef{7}）。
\end{minipage}

\bksubsection{8}{平均の値}

\label{bk:5-12:heikin}自転車が、0秒から2秒まで $v=3t^2$（m/s）で速くなりました。この2秒の平均の速さはいくつでしょう。

\noindent\textbf{方法1　たくさんの値の平均}\quad 2-6の平均は「合計÷個数」でした。2秒を4つに分け、それぞれのまん中の時刻（0.25、0.75、1.25、1.75秒）の速さをはかると、0.1875、1.6875、4.6875、9.1875です。平均は
\[ (0.1875+1.6875+4.6875+9.1875)\div4=15.75\div4=3.9375 \]
8つに分けると3.984…、もっと細かくすると4に近づきます。

\noindent\textbf{方法2　面積を長方形にならす}\quad 進んだ距離は $\displaystyle\int_0^2 3t^2\,\dd t=\Bigl[\,t^3\,\Bigr]_0^2=8$（m）です。2秒で\qty{8}{m}進んだので、平均の速さは $8\div2=\qty{4}{m/s}$。5-1の「全部の距離÷全部の時間」です。

\noindent\begin{minipage}[c]{0.53\linewidth}\centering
\begin{tikzpicture}
\begin{axis}[bkplot, width=60mm, height=38mm, xmin=0, xmax=2.4, ymin=0, ymax=13, xtick={1,2}, ytick={4,12}, xlabel={$t$（s）}, ylabel={$v$（m/s）}]
  \fill[bkOrange!25] (axis cs:0,0) rectangle (axis cs:2,4);
  \addplot[fill=bkBlue!20, fill opacity=0.7, draw=none, domain=0:2] {3*x^2} \closedcycle;
  \addplot[bkcurve, domain=0:2.08] {3*x^2};
  \draw[bkOrange, line width=1.2pt] (axis cs:0,4) -- (axis cs:2,4);
  \node[font=\sffamily\scriptsize, anchor=west, text=bkOrange!80!black] at (axis cs:2.08,4) {平均 4};
\end{axis}
\end{tikzpicture}
\bkfigcaption{青とオレンジの面積は、どちらも8。}
\end{minipage}\hfill
\begin{minipage}[c]{0.44\linewidth}\raggedright
グラフでは、平均の値は、もとのグラフの下の面積と同じ面積をもつ長方形の高さです。でこぼこの砂をならして、平らにした高さと同じです。
\end{minipage}

\noindent\textbf{2つの方法がつながるわけ}\quad 方法1で、はば $\Delta t$ の区間に $n$ こ分けたとき、値の合計÷個数の分子と分母に $\Delta t$ をかけると（分数の大きさは変わらない）、
\[ \frac{\text{値の合計}}{n}=\frac{\text{値の合計}\times\Delta t}{n\times\Delta t}=\frac{\text{細い長方形の面積の和}}{\text{全体のはば}} \]
になります。はばを0に近づけると、分子は積分になります。だから
\[ \text{平均の値}=\frac{1}{b-a}\int_a^b f(t)\,\dd t \]
です。$b-a$ でわるのは、「個数でわる」のかわりに「全体のはば（時間）でわる」からです。

\begin{bkcenter}
\begin{bktable}{colspec={Q[l,m,wd=36mm] Q[l,m,wd=36mm] Q[c,m,wd=36mm]}, rowsep=1pt, row{1}={bg=bkTheme!10}}
1あたりの量 & 積分すると（たまる量） & 単位の計算 \\
速度 $v$（m/s） & 位置の変化（m） & (m/s)×s＝m \\
力 $F$（N） & 仕事 $W$（J） & N×m＝J \\
流量 $Q$（L/min） & 水の量 $V$（L） & (L/min)×min＝L \\
電力 $P$（W） & 電力量 $E$（J） & (J/s)×s＝J \\
電流 $I$（A） & 電気の量 $q$（C） & (C/s)×s＝C \\
\end{bktable}
\bktabcaption{積分は「1あたりの量」から「たまる量」を求める道具。}
\end{bkcenter}

単位の計算では、かけ算で「あたり」（/s や /min）が消えます。力だけは「1あたり」に見えませんが、仕事を位置で細かく分けると「\qty{1}{m}動かすあたりの仕事」が力になっていて、同じ形です。

\begin{column}{スマートフォンの電池の「mAh」}
電池の箱の「\qty{3000}{mAh}（ミリアンペア時）」は、電流（mA）を時間（h）で積分した、とり出せる電気の量です。\qty{1}{mA}の電流が1時間流れた電気の量が\qty{1}{mAh}です。\qty{300}{mA}の電流で使い続けると、
\[ \qty{3000}{mAh}\div\qty{300}{mA}=\qty{10}{h} \]
で、10時間使えます。電流が変わっても、積分した値が\qty{3000}{mAh}になるまで使えます。\qty{1}{mAh}を C に直すと、$\qty{0.001}{A}\times\qty{3600}{s}=\qty{3.6}{C}$ です。
\end{column}

\noindent\textbf{3つの段のまとめ}\quad この章で求めた量は、どれも同じ3つの段で求めています。ちがうのは、何で細かく分けるか（位置か時間か）と、何をかけるかだけです。
\begin{bkcenter}
\begin{bktable}{colspec={Q[l,m,wd=22mm] Q[l,m,wd=30mm] Q[c,m,wd=30mm] Q[c,m,wd=26mm]}, rowsep=2pt, row{1}={bg=bkTheme!10}}
求める量 & 段1　分ける & 段2　かける & 段3　たす \\
仕事 & 位置を細かく（$\dd x$） & 力×小さな距離 $F\,\dd x$ & $\displaystyle\int F\,\dd x$ \\
水の量 & 時間を細かく（$\dd t$） & 流量×小さな時間 $Q\,\dd t$ & $\displaystyle\int Q\,\dd t$ \\
電力量 & 時間を細かく（$\dd t$） & 電力×小さな時間 $P\,\dd t$ & $\displaystyle\int P\,\dd t$ \\
電気の量 & 時間を細かく（$\dd t$） & 電流×小さな時間 $I\,\dd t$ & $\displaystyle\int I\,\dd t$ \\
距離（5-10） & 時間を細かく（$\dd t$） & 速度×小さな時間 $v\,\dd t$ & $\displaystyle\int v\,\dd t$ \\
\end{bktable}
\bktabcaption{どの量も「分ける・かける・たす」。たした先が定積分。}
\end{bkcenter}

\begin{taisetsu}
変わる量の積分 ＝ たまる量（細かく分けて、かけて、たす）\\
\quad 仕事 $\displaystyle\int F\,\dd x$、水の量 $\displaystyle\int Q\,\dd t$、電力量 $\displaystyle\int P\,\dd t$、電気の量 $\displaystyle\int I\,\dd t$\\
ばねの仕事 $\dfrac{1}{2}kx^2$\qquad 平均の値 $=\dfrac{1}{b-a}\displaystyle\int_a^b f\,\dd t$
\end{taisetsu}

\begin{kotoba}
\small
\bkword{仕事}{力×力の向きに動いた距離。力が変わるときは $\int F\,\dd x$。単位は J。}
\bkword{流量・電力・電流}{1分（1秒）あたりの水の量・エネルギー・電気の量。積分すると、たまった量・使った量。}
\bkword{電力量}{電力を時間で積分したもの。kWh や J で表す。}
\end{kotoba}

%% ---------------- 例題 ----------------
\bksection{reidai}{例題でたしかめよう}

\begin{reidai}[一定の力の仕事]
質量\qty{5}{kg}の荷物を、ゆっくり一定の速さで\qty{2}{m}の高さまで持ち上げる仕事を求めましょう（重力は1 kg あたり\qty{9.8}{N}）。\par
\kangaekata 持ち上げる力は荷物の重力と同じで、一定です。だから①の「力×距離」を使います。\par
\shiki 力\quad $F=\qty{5}{kg}\times\qty{9.8}{N/kg}=\qty{49}{N}$\par
\hspace*{2.6\zw}仕事\quad $W=Fx=49\times2=98$\par
\kotae \qty{98}{J}
\end{reidai}

\begin{jibun}
\qty{30}{N}の力で、台車を力の向きに\qty{4}{m}おしました。仕事を求めましょう。
\bkjans{$30\times4=\qty{120}{J}$}
\end{jibun}

\noindent\begin{minipage}[c]{0.46\linewidth}
\begin{tikzpicture}
\begin{axis}[bkplot, width=52mm, height=40mm, xmin=0, xmax=2.6, ymin=0, ymax=62, xtick={1,2}, ytick={49}, xlabel={$x$（m）}, ylabel={$F$（N）}]
  \fill[bkOrange!25] (axis cs:0,0) rectangle (axis cs:2,49);
  \addplot[bkcurve, line width=1.4pt] coordinates {(0,49) (2,49)};
  \node[font=\sffamily\small, fill=bkOrange!25, inner sep=1pt] at (axis cs:1,22) {\qty{98}{J}};
\end{axis}
\end{tikzpicture}
\end{minipage}\hfill
\begin{minipage}[c]{0.52\linewidth}\raggedright
\bkreidairef{1}を $F$-$x$ グラフで見ると、たて\qty{49}{N}、横\qty{2}{m}の長方形です。面積 $49\times2=98$ が仕事です。力が一定のときは、細かく分けなくても、長方形1つで求められます。
\end{minipage}

\begin{reidai}[ばねを、とちゅうからのばす仕事]
ばね定数 $k=\qty{200}{N/m}$ のばねを、のび\qty{0.1}{m}から\qty{0.2}{m}までのばす仕事を求めましょう。\par
\kangaekata 力 $F=200x$ はのびとともに変わるので、③の $\int F\,\dd x$ を使います。はじめの位置0.1が下端、おわりの位置0.2が上端です。\par
\shiki $\displaystyle\int_{0.1}^{0.2}200x\,\dd x=\Bigl[\,200\times\frac{x^2}{2}\,\Bigr]_{0.1}^{0.2}=\Bigl[\,100x^2\,\Bigr]_{0.1}^{0.2}$\par
\hspace*{2.6\zw}$=100\times0.2^2-100\times0.1^2=100\times0.04-100\times0.01=4-1=3$\par
\hspace*{2.6\zw}確かめ：台形の面積（力は20 Nから40 N）\par
\hspace*{2.6\zw}$(20+40)\times0.1\div2=3$\par
\kotae \qty{3}{J}（0から\qty{0.1}{m}のときの3倍。のびているほど力が大きいため）
\end{reidai}

\begin{jibun}
ばね定数 $k=\qty{500}{N/m}$ のばねを、のび0から\qty{0.2}{m}までのばす仕事を求めましょう。
\bkjans{$\dfrac{1}{2}\times500\times0.2^2=\dfrac{1}{2}\times500\times0.04=\qty{10}{J}$}
\end{jibun}

\begin{reidai}[グラフが曲がっている力の仕事]
位置 $x$（m）で力が $F=3x^2$（N）になる物体を、$x=0$ から2まで動かす仕事を求めましょう。\par
\kangaekata $F$-$x$ グラフは曲がっているので、三角形や台形では求められません。③の $\int F\,\dd x$（「力×小さな距離」の和）を、5-11の方法で計算します。\par
\shiki $3x^2$ を積分すると $3\times\dfrac{x^3}{3}=x^3$（指数を1ふやして、その数でわる）\par
\hspace*{2.6\zw}$\displaystyle\int_0^2 3x^2\,\dd x=\Bigl[\,x^3\,\Bigr]_0^2=2^3-0^3=8-0=8$\par
\kotae \qty{8}{J}
\end{reidai}

\begin{jibun}
力が $F=4x+2$（N）のとき、$x=0$ から3まで動かす仕事を求めましょう。
\bkjans{$\Bigl[\,2x^2+2x\,\Bigr]_0^3=18+6=\qty{24}{J}$}
\end{jibun}

\begin{reidai}[流量から水の量]
ポンプでタンクに水を入れます。流量が $Q=3t^2$（L/min、$t$ は分）でふえていきます。0分から4分までに入った水の量を求めましょう。\par
\kangaekata 流量（L/min）を時間（min）で積分すると、水の量。流量が「1分あたり」で、時間も分なので、単位はそろっています。\par
\shiki $\displaystyle\int_0^4 3t^2\,\dd t=\Bigl[\,3\times\frac{t^3}{3}\,\Bigr]_0^4=\Bigl[\,t^3\,\Bigr]_0^4=4^3-0=64$\par
\hspace*{2.6\zw}単位\quad (L/min)×min＝L\par
\kotae \qty{64}{L}
\end{reidai}

\begin{jibun}
大きなポンプの流量が $Q=10-2t$（L/s、$t$ は秒）でへっていき、5秒で止まりました。出た水の量を求めましょう。
\bkjans{$\Bigl[\,10t-t^2\,\Bigr]_0^5=50-25=\qty{25}{L}$}
\end{jibun}

\noindent\begin{minipage}[c]{0.46\linewidth}
\begin{tikzpicture}
\begin{axis}[bkplot, width=54mm, height=38mm, xmin=0, xmax=4.6, ymin=0, ymax=54, xtick={1,2,3,4}, ytick={12,27,48}, xlabel={$t$（min）}, ylabel={$Q$（L/min）}]
  \foreach \k in {1,...,4} {\edef\tmp{\noexpand\fill[bkBlue!15, draw=bkBlue!60, line width=0.4pt] (axis cs:\k-1,0) rectangle (axis cs:\k,{3*\k*\k});}\tmp}
  \addplot[fill=bkBlue!35, draw=none, domain=0:4] {3*x^2} \closedcycle;
  \addplot[bkcurve, domain=0:4.2] {3*x^2};
\end{axis}
\end{tikzpicture}
\end{minipage}\hfill
\begin{minipage}[c]{0.52\linewidth}\raggedright
\bkreidairef{4}を1分ずつ4つに分け、右はしの流量×1分をたすと $3+12+27+48=90$（大きめ）、左はしでは $0+3+12+27=42$（小さめ）。ほんとうの64 L（こい青）は、そのあいだです。
\end{minipage}

\begin{reidai}[電力が一定のときの電力量]
\qty{1.2}{kW}のドライヤーを15分使いました。電力量は何 kWh で、何 J ですか。\par
\kangaekata 電力が一定なので、電力量＝電力×時間（⑥）。kWh で答えるので、時間を h（時間）に直してからかけます。J には⑥の $\qty{1}{kWh}=\qty{3.6e6}{J}$ で直します。\par
\shiki 時間\quad $\qty{15}{min}\div\qty{60}{min/h}=\qty{0.25}{h}$\par
\hspace*{2.6\zw}電力量\quad $\qty{1.2}{kW}\times\qty{0.25}{h}=\qty{0.3}{kWh}$\par
\hspace*{2.6\zw}J に直す\quad $0.3\times\qty{3.6e6}{J}=\qty{1.08e6}{J}$\par
\kotae \qty{0.3}{kWh}、\qty{1.08e6}{J}（108万 J）
\end{reidai}

\begin{jibun}
\qty{800}{W}の電子レンジを3分使いました。電力量は何 kWh で、何 J ですか。
\bkjans{$\qty{0.8}{kW}\times\qty{0.05}{h}=\qty{0.04}{kWh}$、$800\times180=\qty{144000}{J}$}
\end{jibun}

\begin{reidai}[電力量と平均の電力]
電力が $P=6-t$（kW）の充電器で、0時間から4時間まで充電しました。電力量と、平均の電力を求めましょう。\par
\kangaekata 電力を時間で積分すると電力量（⑥）。平均は、⑧の「積分÷時間」です。kW×h なので、単位は kWh。\par
\shiki $\displaystyle\int_0^4(6-t)\,\dd t=\Bigl[\,6t-\frac{t^2}{2}\,\Bigr]_0^4=\Bigl(24-\frac{16}{2}\Bigr)-0=24-8=16$\par
\hspace*{2.6\zw}平均\quad $16\div4=4$\par
\kotae \qty{16}{kWh}、平均\qty{4}{kW}（はじめの6 kWとおわりの2 kWのまん中）
\end{reidai}

\begin{jibun}
電力が $P=2t$（kW）でふえるヒーターを、0時間から3時間まで使いました。電力量と平均の電力を求めましょう。
\bkjans{$\Bigl[\,t^2\,\Bigr]_0^3=\qty{9}{kWh}$、平均 $9\div3=\qty{3}{kW}$}
\end{jibun}

\begin{reidai}[電流から電気の量]
コンデンサー（電気をためる部品）に流れこむ電流が $I=2-0.5t$（A）でへっていき、4秒で0になりました。たまった電気の量 $q$ を求めましょう。\par
\kangaekata 電流（A＝C/s）は1秒あたりの電気の量（⑦）。時間で積分すると電気の量（C）。\par
\shiki $\displaystyle\int_0^4(2-0.5t)\,\dd t=\Bigl[\,2t-0.25t^2\,\Bigr]_0^4$\par
\hspace*{2.6\zw}$=(2\times4-0.25\times16)-0=8-4=4$\par
\hspace*{2.6\zw}確かめ：三角形の面積 $\dfrac{1}{2}\times4\times2=4$\par
\kotae \qty{4}{C}
\end{reidai}

\begin{jibun}
電流が $I=0.3t$（A）でふえていきます。0秒から10秒までに流れた電気の量を求めましょう。
\bkjans{$\Bigl[\,0.15t^2\,\Bigr]_0^{10}=\qty{15}{C}$}
\end{jibun}

\begin{reidai}[変わる量の平均]
信号で止まっていた自転車が走り出し、次の信号で止まりました。速さは $v=6t-t^2$（m/s、0秒から6秒まで）でした。平均の速さを求めましょう。\par
\kangaekata 平均＝積分÷時間（⑧）。はじめと最後の速さはどちらも0なので、「(はじめ＋おわり)÷2」では求められません。\par
\shiki $\displaystyle\int_0^6(6t-t^2)\,\dd t=\Bigl[\,3t^2-\frac{t^3}{3}\,\Bigr]_0^6=\Bigl(3\times36-\frac{216}{3}\Bigr)-0$\par
\hspace*{2.6\zw}$=108-72=36$\par
\hspace*{2.6\zw}平均\quad $36\div6=6$\par
\kotae \qty{6}{m/s}（6秒で\qty{36}{m}進んだ）
\end{reidai}

\begin{jibun}
速度が $v=3t^2$（m/s）の物体の、0秒から3秒までの平均の速さを求めましょう。
\bkjans{$\Bigl[\,t^3\,\Bigr]_0^3=27$。$27\div3=\qty{9}{m/s}$}
\end{jibun}

\noindent\begin{minipage}[c]{0.48\linewidth}
\begin{tikzpicture}
\begin{axis}[bkplot, width=56mm, height=42mm, xmin=0, xmax=7, ymin=0, ymax=11, xtick={1,2,...,6}, ytick={6,9}, xlabel={$t$（s）}, ylabel={$v$（m/s）}]
  \fill[bkOrange!25] (axis cs:0,0) rectangle (axis cs:6,6);
  \addplot[fill=bkBlue!20, fill opacity=0.7, draw=none, domain=0:6] {6*x-x^2} \closedcycle;
  \addplot[bkcurve, domain=0:6] {6*x-x^2};
  \draw[bkOrange, line width=1.2pt] (axis cs:0,6) -- (axis cs:6,6);
\end{axis}
\end{tikzpicture}
\end{minipage}\hfill
\begin{minipage}[c]{0.5\linewidth}\raggedright
\bkreidairef{8}の自転車の $v$-$t$ グラフです。速さは0から\qty{9}{m/s}まで上がって、また0にもどります。青い部分の面積（\qty{36}{m}）を、はば6秒の長方形にならすと、高さは6です（オレンジの線）。これが平均の速さ\qty{6}{m/s}です。
\end{minipage}

\begin{reidai}[電池の容量 mAh]
電流が $I=100+50t$（mA、$t$ は時間）で流れたとき、2時間で使った電気の量は、何 mAh（ミリアンペア時）でしょう。\par
\kangaekata 電流（mA）を時間（h）で積分すると、電気の量（mAh）。\qty{1}{mA}が1時間流れた量が\qty{1}{mAh}です。\par
\shiki $\displaystyle\int_0^2(100+50t)\,\dd t=\Bigl[\,100t+50\times\frac{t^2}{2}\,\Bigr]_0^2=\Bigl[\,100t+25t^2\,\Bigr]_0^2$\par
\hspace*{2.6\zw}$=(100\times2+25\times4)-0=200+100=300$\par
\hspace*{2.6\zw}確かめ：台形の面積（電流は100 mAから200 mA）\par
\hspace*{2.6\zw}$(100+200)\times2\div2=300$\par
\kotae \qty{300}{mAh}
\end{reidai}

\begin{jibun}
電流が $I=400-100t$（mA）でへっていき、4時間で0になりました。使った電気の量を求めましょう。
\bkjans{$\Bigl[\,400t-50t^2\,\Bigr]_0^4=1600-800=\qty{800}{mAh}$}
\end{jibun}

\noindent\begin{minipage}[c]{0.46\linewidth}
\begin{tikzpicture}
\begin{axis}[bkplot, width=54mm, height=36mm, xmin=0, xmax=2.7, ymin=0, ymax=260, xtick={1,2}, ytick={100,200}, xlabel={$t$（h）}, ylabel={$I$（mA）}]
  \addplot[fill=bkGreen!25, draw=none, domain=0:2] {100+50*x} \closedcycle;
  \addplot[bkcurve, domain=0:2.5] {100+50*x};
  \draw[bkGray, dashed, line width=0.5pt] (axis cs:2,0) -- (axis cs:2,200);
  \node[font=\sffamily\small, fill=bkGreen!25, inner sep=1pt] at (axis cs:1,70) {\qty{300}{mAh}};
\end{axis}
\end{tikzpicture}
\end{minipage}\hfill
\begin{minipage}[c]{0.5\linewidth}\raggedright
\bkreidairef{9}の $I$-$t$ グラフです。緑の台形の面積が、使った電気の量\qty{300}{mAh}です。たての単位 mA と横の単位 h をかけるので、面積の単位が mAh になります。
\end{minipage}

\begin{reidai}[はかった値から、たまった量を求める]
ポンプの流量を、1分ごとにはかると、下の表のようでした。0分から4分までに出た水の量を、およそ求めましょう。
\begin{bkcenter}
\begin{bktable}{colspec={Q[l,m,wd=30mm] *{5}{Q[c,m,wd=10mm]}}, rowsep=1pt}
$t$（min） & 0 & 1 & 2 & 3 & 4 \\
$Q$（L/min） & 2 & 6 & 8 & 9 & 9 \\
\end{bktable}
\end{bkcenter}
\kangaekata 式がないので、積分の公式は使えません。そこで、5-10のように、細い図形の面積をたします。1分ごとの区間の中では、流量はほぼまっすぐ変わると考え、はじめとおわりの流量を上底・下底とする\textbf{台形}の面積を、その1分の水の量とします。台形の面積は、左はしの長方形（小さめ）と右はしの長方形（大きめ）の、ちょうどまん中です。\par
\shiki 0〜1分\quad $(2+6)\times1\div2=4$\par
\hspace*{2.6\zw}1〜2分\quad $(6+8)\times1\div2=7$\par
\hspace*{2.6\zw}2〜3分\quad $(8+9)\times1\div2=8.5$\par
\hspace*{2.6\zw}3〜4分\quad $(9+9)\times1\div2=9$\par
\hspace*{2.6\zw}合計\quad $4+7+8.5+9=28.5$\par
\hspace*{2.6\zw}確かめ：右はしの和 $6+8+9+9=32$\par
\hspace*{2.6\zw}\phantom{確かめ：}左はしの和 $2+6+8+9=25$。28.5はそのあいだ\par
\kotae 約\qty{28.5}{L}
\end{reidai}

\begin{jibun}
太陽光パネルの電力を1時間ごとにはかると、0、2、3、3、\qty{2}{kW}（0時間から4時間まで）でした。4時間の電力量を、台形の面積の和でおよそ求めましょう。
\bkjans{$1+2.5+3+2.5=\qty{9}{kWh}$}
\end{jibun}

\noindent\begin{minipage}[c]{0.46\linewidth}
\begin{tikzpicture}
\begin{axis}[bkplot, width=54mm, height=38mm, xmin=0, xmax=4.6, ymin=0, ymax=11, xtick={1,2,3,4}, ytick={2,6,9}, xlabel={$t$（min）}, ylabel={$Q$（L/min）}]
  \addplot[fill=bkGreen!25, draw=bkGreen!60!black, line width=0.4pt] coordinates {(0,0) (0,2) (1,6) (2,8) (3,9) (4,9) (4,0)} -- cycle;
  \draw[bkGreen!60!black, line width=0.4pt] (axis cs:1,0) -- (axis cs:1,6) (axis cs:2,0) -- (axis cs:2,8) (axis cs:3,0) -- (axis cs:3,9);
  \addplot[only marks, mark=*, mark size=1.6pt, bkInk] coordinates {(0,2) (1,6) (2,8) (3,9) (4,9)};
\end{axis}
\end{tikzpicture}
\end{minipage}\hfill
\begin{minipage}[c]{0.5\linewidth}\raggedright
\bkreidairef{10}の表の点を直線でつなぐと、4つの台形（緑）ができます。その面積の和が28.5 L です。工場の積算流量計や家の電力量計（スマートメーター）も、はかった値を、ごく短い時間ごとにこのようにたし合わせています。
\end{minipage}

\begin{reidai}[階段をのぼる仕事]
質量\qty{40}{kg}の人が、高さ\qty{3}{m}（建物の1階分）の階段を、ゆっくり一定の速さでのぼります。重力にさからってする仕事は何 J ですか。また、それは何 kWh ですか（重力は1 kg あたり\qty{9.8}{N}）。\par
\kangaekata 階段はななめに進みますが、体を持ち上げる力は上向きです。力の向き（上向き）に動いた距離は、のぼった高さの\qty{3}{m}です。横に進んだ分は、上向きの力の向きに動いていないので、この仕事に入りません。kWh には、⑥の $\qty{1}{kWh}=\qty{3.6e6}{J}$ でわって直します。\par
\shiki 力\quad $F=\qty{40}{kg}\times\qty{9.8}{N/kg}=\qty{392}{N}$\par
\hspace*{2.6\zw}仕事\quad $W=392\times3=\qty{1176}{J}$\par
\hspace*{2.6\zw}kWh に直す\quad $\qty{1176}{J}\div(\qty{3.6e6}{J/kWh})\approx0.00033$\,kWh\par
\kotae \qty{1176}{J}（約\qty{1.2}{kJ}）、約0.00033\,kWh
\end{reidai}

\noindent\begin{minipage}[c]{0.44\linewidth}\centering
\begin{tikzpicture}[x=1mm, y=1mm]
  \fill[bkGray!25] (0,0) -- (0,6) -- (6,6) -- (6,12) -- (12,12) -- (12,18) -- (18,18) -- (18,24) -- (24,24) -- (24,30) -- (32,30) -- (32,0) -- cycle;
  \draw[bkInk, line width=0.6pt] (-2,0) -- (0,0) -- (0,6) -- (6,6) -- (6,12) -- (12,12) -- (12,18) -- (18,18) -- (18,24) -- (24,24) -- (24,30) -- (32,30);
  \bkperson{(2,6.5)}{13}{bkBlue}
  \draw[bk arrow, bkOrange, line width=0.9pt] (-5,4) -- (-5,30);
  \node[font=\sffamily\scriptsize, anchor=south, text=bkOrange!80!black] at (-5,31) {力（上向き）};
  \draw[bkGray, <->, line width=0.5pt] (35,0) -- (35,30);
  \node[font=\sffamily\scriptsize, anchor=west] at (36.5,15) {\qty{3}{m}};
\end{tikzpicture}
\end{minipage}\hfill
\begin{minipage}[c]{0.53\linewidth}\raggedright
\bkreidairef{11}の仕事1176 J は、電気代にすると約0.01円（1 kWh を31円として、$0.00033\times31\approx0.01$）。1 kWh は、この階段をおよそ3000回のぼる仕事です（$\qty{3.6e6}{J}\div\qty{1176}{J}\approx3061$）。kWh がとても大きな単位だとわかります。
\end{minipage}

\begin{jibun}
質量\qty{50}{kg}の人が、高さ\qty{12}{m}（4階分）の階段をのぼります。重力にさからってする仕事は何 J で、何 kWh ですか（重力は1 kg あたり\qty{9.8}{N}）。
\bkjans{$50\times9.8\times12=\qty{5880}{J}$、$5880\div(3.6\times10^6)\approx\qty{0.0016}{kWh}$}
\end{jibun}

%% ---------------- まちがえやすいところ・物理工学 ----------------
\bksection{tsumazuki}{まちがえやすいところをチェック}

\begin{tsumazuki}{ばねの仕事を $kx\times x$ としてしまう}
\bkng{$W=kx\times x=kx^2$}\bkok{$W=\dfrac{1}{2}kx^2$}
$kx$ は、のびが $x$ になったときの最後の力で、はじめの力は0です。$F$-$x$ グラフは三角形なので、$\dfrac{1}{2}$ がつきます。
\end{tsumazuki}

\begin{tsumazuki}{時間の単位と流量の単位がそろっていない}
\bkng{$Q=\qty{12}{L/min}$ で10秒なので $12\times10=\qty{120}{L}$}\bkok{$\qty{12}{L/min}\times\qty{10}{s}\div\qty{60}{s/min}=\qty{2}{L}$}
流量が「1分あたり」なのに秒をかけると、60倍まちがえます。単位を書いてかければ、min と s が約分できないことに気づきます（2-10）。
\end{tsumazuki}

\begin{tsumazuki}{kW と kWh を同じものと考える}
\bkng{\qty{1.2}{kW}のドライヤーを15分使ったので、電力量は\qty{1.2}{kWh}}\bkok{$\qty{1.2}{kW}\times\qty{0.25}{h}=\qty{0.3}{kWh}$}
kW は電力（1秒あたりに使うエネルギー）で、速さのような「あたり」の量です。kWh は電力×時間で、進んだ距離のような「たまった量」です。時間をかけないと kWh にはなりません。
\end{tsumazuki}

\begin{tsumazuki}{平均を「(はじめ＋おわり)÷2」としてしまう}
\bkng{$v=3t^2$ の0〜2秒の平均は $(0+12)\div2=\qty{6}{m/s}$}\bkok{$8\div2=\qty{4}{m/s}$（積分÷時間）}
「(はじめ＋おわり)÷2」が使えるのは、グラフがまっすぐなとき（台形）だけです。曲がっているときは、積分÷時間で求めます。
\end{tsumazuki}

\begin{column}{電気代を計算しよう}
電力会社は、使った電力量（kWh）に応じて電気代を決めます。1 kWh あたりの料金を約31円とします。\qty{1.2}{kW}のドライヤーを毎日15分（\qty{0.25}{h}）、30日使うと、
\[ \qty{1.2}{kW}\times\qty{0.25}{h}\times30=\qty{9}{kWh},\qquad 9\times31=279 \]
で、1か月の電気代は約280円です。電力（kW）が大きい器具ほど、また長い時間使うほど、電力量がたまって電気代が高くなります。
\end{column}

\begin{butsuri}
\begin{itemize}[leftmargin=1.2\zw, itemsep=1pt, topsep=0pt, label={\color{bkPurple}●}]
\item \textbf{電気のメーター}\\家の電力量計（スマートメーター）は、電力を細かい時間ごとにはかってたし合わせ（積分し）、kWh で表示しています。
\item \textbf{ばねのエネルギー}\\おもちゃの車のぜんまいや、時計のばね、車のサスペンションは、$\frac{1}{2}kx^2$ のエネルギーをたくわえたり、とり出したりしています。
\item \textbf{工場の積算流量計}\\化学工場では、パイプを流れる液体やガスの流量を時間で積分して、原料や製品の量を管理しています。
\end{itemize}
\end{butsuri}

%% ---------------- 章末問題 ----------------
\bkendproblems
\bkproblemgroup{基本}{3}

\bkproblem \qty{20}{N}の力で、物を力の向きに\qty{3}{m}動かしました。仕事を求めましょう。
\bkpans{\qty{60}{J}}

\bkproblem ばね定数 $k=\qty{400}{N/m}$ のばねについて、（1）0から\qty{0.05}{m}、（2）\qty{0.05}{m}から\qty{0.1}{m} のばす仕事を求めましょう。
\bkpans{（1）\nobreak \qty{0.5}{J}\quad（2）\nobreak $2-0.5=\qty{1.5}{J}$}

\bkproblem ポンプの流量が $Q=0.5t$（L/s）でふえていきます。0秒から20秒までに出た水の量を求めましょう。
\bkpans{$\Bigl[\,0.25t^2\,\Bigr]_0^{20}=\qty{100}{L}$}

\bkproblem \qty{1.2}{kW}のドライヤーを10分使った電力量は、何 kWh、何 J ですか。
\bkpans{$1.2\times\dfrac{10}{60}=\qty{0.2}{kWh}$、$0.2\times3.6\times10^6=\qty{7.2e5}{J}$}

\bkproblem 位置 $x$（m）で力が $F=6x$（N）になる物体を、$x=0$ から2まで動かす仕事を求めましょう。
\bkpans{$\Bigl[\,3x^2\,\Bigr]_0^2=\qty{12}{J}$}

\bkproblem 電流が $I=0.2t$（A）でふえていきます。0秒から5秒までに流れた電気の量を求めましょう。
\bkpans{$\Bigl[\,0.1t^2\,\Bigr]_0^5=\qty{2.5}{C}$}

\bkproblemgroup{チャレンジ}{3}

\bkproblem 速度が $v=t^2$（m/s）の物体の、0秒から3秒までの平均の速さを求めましょう。
\bkpans{距離 $\Bigl[\,\dfrac{t^3}{3}\,\Bigr]_0^3=9$。平均 $9\div3=\qty{3}{m/s}$}

\bkproblem 小さな太陽光パネルの電力が、朝6時から夕方6時まで $P=12t-t^2$（W、$t$ は6時からの時間）で変わりました。1日の電力量（Wh）と、平均の電力を求めましょう。
\bkpans{$\Bigl[\,6t^2-\dfrac{t^3}{3}\,\Bigr]_0^{12}=864-576=\qty{288}{Wh}$、平均 $288\div12=\qty{24}{W}$}

\bkproblem スマートフォンから、電流 $I=300+100t$（mA、$t$ は時間）が流れ続けたとします。4時間で使った電気の量は何 mAh ですか。\qty{3000}{mAh}の電池は、4時間もちますか。
\bkpans{$\Bigl[\,300t+50t^2\,\Bigr]_0^4=1200+800=\qty{2000}{mAh}$。3000より少ないので、\mbox{4時間}もつ}

\bkproblem 暖房の強さを変えた部屋の温度が $T=15+6t-t^2$（℃、$t$ は時間）で変わりました。$t=0$ から6までの平均の温度を求めましょう。
\bkpans{$\Bigl[\,15t+3t^2-\dfrac{t^3}{3}\,\Bigr]_0^6=90+108-72=126$。$126\div6=\mbox{$21$\,℃}$}

\begin{furikaeri}
\begin{checklist}
\item 変わる力の仕事を、積分で求められる。\hfill{\footnotesize\sffamily\color{bkTheme}→\pageref{bk:5-12:sekibun}ページ}
\item 流量や電力から、たまった量を求められる。\hfill{\footnotesize\sffamily\color{bkTheme}→\pageref{bk:5-12:ryuuryou}ページ}
\item 積分を使って、平均の値を求められる。\hfill{\footnotesize\sffamily\color{bkTheme}→\pageref{bk:5-12:heikin}ページ}
\end{checklist}
\end{furikaeri}

\bkfukidashi{mina}{たまった量は積分でわかるけど、「はじめにどれだけあったか」がわからないと、いまの量は決まらないよね。}
